% ----------------------------------------------------------
% Metodologia
% ----------------------------------------------------------
\chapter{METODOLOGIA}

A metodologia deste trabalho foi organizada de acordo com as etapas necessárias para compreender o problema, definir as funcionalidades, desenvolver o protótipo e verificar seu funcionamento. 

\section{LEVANTAMENTO BIBLIOGRÁFICO E DEFINIÇÃO DOS REQUISITOS}

A primeira etapa corresponde ao levantamento bibliográfico, realizado a partir de artigos, trabalhos acadêmicos e relatórios relacionados a jogos digitais, conteúdos produzidos por jogadores, organização de informações e gestão de dados em aplicações web. 

Com base nos estudos analisados, serão definidos os requisitos funcionais e não funcionais do protótipo. Entre as principais funcionalidades estão o cadastro de usuários e jogos, a publicação de contribuições vinculadas a cada título, a classificação por categorias e a recuperação das informações por meio de consultas e filtros.

\section{MODELAGEM E DESENVOLVIMENTO DO PROTÓTIPO}

Após a definição dos requisitos, serão elaborados a arquitetura, o modelo de dados e os principais fluxos de funcionamento. O modelo será organizado de forma que cada publicação permaneça vinculada ao jogador responsável, ao jogo correspondente e à categoria selecionada, como sugestão, ideia, avaliação ou relato de problema.

O protótipo será desenvolvido utilizando TypeScript no ambiente Node.js, com o framework NestJS para a organização das funcionalidades. O PostgreSQL será utilizado para o armazenamento dos dados, enquanto o Prisma será responsável pela comunicação com o banco e pelo gerenciamento de sua estrutura. As operações disponibilizadas serão documentadas por meio do Swagger.

A Figura \ref{fig:metodologia} sintetiza as etapas metodológicas adotadas para o desenvolvimento do protótipo.

\begin{figure}[htbp]
\centering
\resizebox{\textwidth}{!}{%
\begin{tikzpicture}[
    font=\footnotesize,
    bloco/.style={
        draw,
        rounded corners=3pt,
        align=center,
        text width=4.2cm,
        minimum height=2.3cm,
        fill=white,
        inner sep=6pt,
        line width=0.6pt
    },
    etapa/.style={
        draw,
        rounded corners=4pt,
        fill=gray!8,
        line width=0.7pt
    },
    titulo/.style={
        font=\bfseries\small,
        align=center
    },
    seta/.style={
        -{Latex[length=2.5mm,width=1.8mm]},
        line width=0.7pt
    }
]

% Fundos das etapas (desenhados primeiro para permanecerem atrás do conteúdo)
\draw[etapa] (-7.4,0.5) rectangle (7.4,-4.8);
\draw[etapa] (-7.4,-5.0) rectangle (7.4,-9.45);

% Etapa 1
\node[titulo] at (0,0) {Etapa 1: Levantamento bibliográfico e definição dos requisitos};

\node[bloco] (levantamento) at (-2.7,-2.45) {
    \textbf{Levantamento bibliográfico}\\[0.15cm]
    Jogos digitais\\
    Conteúdos produzidos por jogadores\\
    Organização de informações\\
    Gestão de dados em aplicações web
};

\node[bloco] (requisitos) at (2.7,-2.45) {
    \textbf{Definição dos requisitos}\\[0.15cm]
    Requisitos funcionais\\
    Requisitos não funcionais\\
    Cadastro de usuários e jogos\\
    Publicação, classificação e consulta
};

% Etapa 2
\node[titulo] at (0,-5.5) {Etapa 2: Modelagem e desenvolvimento do protótipo};

\node[bloco] (modelagem) at (-4.9,-7.6) {
    \textbf{Modelagem da solução}\\[0.15cm]
    Arquitetura da API\\
    Modelo de dados\\
    Fluxos de funcionamento
};

\node[bloco] (organizacao) at (0,-7.6) {
    \textbf{Organização das publicações}\\[0.15cm]
    Jogador responsável\\
    Jogo correspondente\\
    Categoria selecionada\\
    Sugestão, ideia, avaliação ou relato
};

\node[bloco] (desenvolvimento) at (4.9,-7.6) {
    \textbf{Desenvolvimento do protótipo}\\[0.15cm]
    TypeScript e Node.js\\
    NestJS\\
    PostgreSQL e Prisma\\
    Swagger
};

% Setas em espaços livres, sem atravessar textos
\draw[seta] (levantamento.east) -- (requisitos.west);
\draw[seta] (modelagem.east) -- (organizacao.west);
\draw[seta] (organizacao.east) -- (desenvolvimento.west);
\draw[seta] (0,-4.8) -- (0,-5.0);

\end{tikzpicture}
}
\caption{Etapas metodológicas adotadas para o desenvolvimento do protótipo}
\label{fig:metodologia}
\fonte{Elaborado pelo autor (2026).}
\end{figure}
